\documentclass[11pt,a4paper]{article}
\usepackage[utf8]{inputenc}
\usepackage[T1]{fontenc}
\usepackage[french]{babel}
\usepackage{amsmath,amssymb,amsthm}
\usepackage{geometry}
\geometry{margin=2.5cm}

\title{Exercice de Probabilités - Classe de Seconde}
\author{NeoHarness AI}
\date{}

\begin{document}

\maketitle

\section{Rappels de cours}
Une expérience aléatoire est une expérience dont on connaît les résultats possibles (issues) sans pouvoir prédire avec certitude lequel va se produire. L'ensemble de toutes les issues forme l'univers $\Omega$.
\begin{itemize}
    \item \textbf{Équiprobabilité :} $P(A) = \frac{\text{nombre d'issues favorables à } A}{\text{nombre total d'issues de } \Omega}$.
    \item \textbf{Propriétés :} $0 \le P(A) \le 1$, $P(\bar{A}) = 1 - P(A)$.
    \item \textbf{Union :} $P(A \cup B) = P(A) + P(B) - P(A \cap B)$.
\end{itemize}

\section{Exercice type et corrigé}

\subsection{Énoncé}
Dans un lycée, on interroge un groupe de 200 élèves sur leur pratique sportive : le football ($F$) et le basketball ($B$).
\begin{itemize}
    \item 80 élèves pratiquent le football.
    \item 50 élèves pratiquent le basketball.
    \item 20 élèves pratiquent les deux sports.
\end{itemize}
On choisit au hasard un élève de ce groupe.

\subsection{Corrigé}
\begin{enumerate}
    \item \textbf{Traduction des données :}
    $$P(F) = \frac{80}{200} = 0,4 \quad ; \quad P(B) = \frac{50}{200} = 0,25 \quad ; \quad P(F \cap B) = \frac{20}{200} = 0,1$$
    \item \textbf{Probabilité de l'union :}
    $$P(F \cup B) = P(F) + P(B) - P(F \cap B) = 0,4 + 0,25 - 0,1 = 0,55$$
    \item \textbf{Probabilité de l'événement contraire :}
    $$P(\overline{F \cup B}) = 1 - P(F \cup B) = 1 - 0,55 = 0,45$$
    \item \textbf{Probabilité de pratiquer uniquement le football :}
    $$\frac{80 - 20}{200} = \frac{60}{200} = 0,3$$
\end{enumerate}

\end{document}
